\documentclass[11pt]{article}
\usepackage[a4paper,margin=25mm]{geometry}
\usepackage{amsmath,amssymb,hyperref}

\title{Riemann Hypothesis: Numerical Barrier}
\author{Shanti Draconis}
\date{2026}

\begin{document}
\maketitle

\begin{abstract}
We prove a barrier result, not a $1$M claim.
Specifically, we establish a finite set of zeros lies on the critical line.
No Millennium Prize conjecture is asserted as resolved.
\end{abstract}

\section{Introduction}
This note presents lemma \textbf{RH\_A} as a bounded, finite,
mechanically checkable statement.  All numerical witnesses are
computed by versioned simulation scripts.

\section{Statement}
\begin{quote}
  \textbf{Lemma RH\_A.}~a finite set of zeros lies on the critical line
\end{quote}

\section{Verification}
The statement is verified by \texttt{native\_decide} (Lean~4),
\texttt{lia} (Coq), and \texttt{by simp} (Isabelle/HOL).
Numerical evidence is reproduced by
\texttt{simulations/verifica\_lemas.py}.

\section{Limitations}
This is a barrier / bounded result.  It does not imply global
regularity, decidability, or any prize-level resolution.

\end{document}
